\documentclass[11pt]{article}
\usepackage[a4paper,margin=23mm]{geometry}
\usepackage[T1]{fontenc}
\usepackage{lmodern,amsmath,amssymb,amsthm,microtype}
\usepackage[colorlinks=true,linkcolor=blue,citecolor=blue,urlcolor=blue]{hyperref}
\usepackage{array,booktabs,tabularx}
\newtheorem{theorem}{Theorem}
\newtheorem{corollary}[theorem]{Corollary}
\newcommand{\C}{\mathbb C}
\setlength{\parskip}{3pt}
\setlength{\emergencystretch}{2em}
\hypersetup{pdftitle={Examples answering Bergman's Question 15 by extending a partition of von zur Gathen and Roche},
pdfauthor={Haobo Ma (ChronoAI Pte Ltd, Singapore); Wenlin Zhang (National University of Singapore; ChronoAI Pte Ltd, Singapore)},
pdfsubject={Licensed under Creative Commons Attribution 4.0 International (CC BY 4.0; cc-by-4.0)}}
\begin{document}
\begin{center}
{\Large\bfseries Examples answering Bergman's Question 15 by extending
a partition of von zur Gathen and Roche}\par
\vspace{9pt}
Haobo Ma\\
{\small ChronoAI Pte Ltd, Singapore}\\[4pt]
Wenlin Zhang\\
{\small National University of Singapore; ChronoAI Pte Ltd, Singapore}\\[6pt]
19 November 2026
\end{center}
\begin{abstract}
We answer Question 15 of George M.\ Bergman's \emph{Thoughts on Eggert's
Conjecture}, arXiv:1206.0326v2, Section 4, printed page 6, in the negative,
including over $\mathbb Q$ by extending a binomial partition of
von zur Gathen and Roche (1997, Table 2, p.\ 359) to a semigroup.
We construct a commutative semigroup with zero having $59$
elements, generated by three elements which have distinct nonzero sixth
powers, while their complex span in the contracted semigroup algebra contains a nonzero
element with sixth power zero. A $49$-element quotient has the same property.
For these two presentations the cancellation exists over a characteristic-zero
field exactly when the field contains a primitive cube root of unity.
The construction also gives examples for every positive exponent divisible by six.
The older partition yields a $632$-element example with two generators and
exponent $35$ over every characteristic-zero field.
\end{abstract}

\section{The question}
Bergman's Question 15~\cite[Section 4, p.\ 6]{Bergman} concerns a finite subset $X$
of a commutative semigroup $S$ with a zero element. The contracted semigroup
algebra $F_0S$, over a field $F$ of characteristic zero, is the vector space
with basis indexed by $S\setminus\{0\}$ and multiplication induced by $S$.
Suppose a positive integer $n$ makes $s\mapsto s^n$ injective on $X$, and no
nonzero element of $X$ has $n$th power zero. Must every nonzero element of
the span $FX\subseteq F_0S$ have nonzero $n$th power?

We use the same symbols for elements of $S$ and their images in $F_0S$.
The map $S\longrightarrow F_0S$ takes $S\setminus\{0\}$ to an $F$-basis
and takes $0$ to the zero vector. Multiplication is extended bilinearly;
a semigroup product equal to zero contributes the zero vector, not another
basis vector. The zero element satisfies $0s=s0=0$. No multiplicative
identity is assumed.

The construction below answers the question in the negative with $F=\C$
and $n=6$. The corresponding locator in the $17$-page published layout
is Section 4, Question 15, p.\ 9.
Section~\ref{sec:older} relates this answer to earlier binomial identities
and gives the rational-field example. For the noncommutative branch of
Question 19 we also specialize the square-zero identity of Almeida,
Margolis, Steinberg and Volkov~\cite[p.\ 11, proof of Proposition 3.4]{AMSV}.
Neither cited paper states these semigroup examples as answers to
Bergman's later questions.

\section{The semigroup}
Let $S$ be the commutative semigroup with zero presented by the generators
\begin{equation}\label{eq:generators}
x,\quad y,\quad z
\end{equation}
and the following relations:
\begin{align}
x^6&=x^5y=x^4y^2=x^3y^3=x^2y^4=xy^5,\label{eq:first}\\
y^6&=y^5z=y^4z^2=y^3z^3=y^2z^4=yz^5,\label{eq:second}\\
z^6&=xz^5=x^2z^4=x^3z^3=x^4z^2=x^5z,\label{eq:third}\\
x^ay^bz^c&=0\quad(a,b,c>0,\ a+b+c=6),\label{eq:mixed}\\
x^ay^bz^c&=0\quad(a,b,c\geq0,\ a+b+c=7).\label{eq:seven}
\end{align}
Every monomial of degree at least seven is consequently zero. We reduce
degree-six expressions by writing $x^6$, $y^6$, $z^6$ for the common values
in~\eqref{eq:first}, \eqref{eq:second}, \eqref{eq:third}, respectively, and
$0$ for the monomials in~\eqref{eq:mixed}.

The semigroup $S$ is obtained from the free commutative semigroup with zero
on $x,y,z$ by setting all monomials of degree greater than six equal to
zero, imposing no relations on monomials of positive degree less than
six, and imposing~\eqref{eq:first}--\eqref{eq:mixed} on the degree-six
monomials. The lists in~\eqref{eq:first}--\eqref{eq:third} are disjoint
and exhaust the degree-six monomials involving at most two generators.
Multiplying a degree-six relation by any positive-degree monomial gives
zero on both sides. The zero class is closed under multiplication, and
the lower-degree singleton classes cause no ambiguity. Hence these
classes form a semigroup congruence. Every monomial of degree at least
seven is a multiple of a degree-seven monomial, so the congruence has
exactly the identifications generated by the displayed relations.
Thus $S$ consists of the distinct monomials of degrees one through five,
the zero element, and the three distinct nonzero degree-six classes
represented by $x^6,y^6,z^6$. As a quotient of a free commutative semigroup
with zero, $S$ is associative and commutative.

There are $\binom{d+2}{2}$ monomials of degree $d$ in three variables, so
\[
|S|=\sum_{d=1}^5\binom{d+2}{2}+3+1=55+3+1=59,
\qquad \dim_{\C}\C_0S=58.
\]
Every product involving $x^6,y^6,z^6$ is zero. Every product of seven
elements is zero, whereas $x^6\ne0$. Thus $S$ and $\C_0S$ have nilpotency
index seven.
For example,
\[
(x^2y)(x^3)=x^6,\quad (y^2z)(y^3)=y^6,\quad
(xyz)(xyz)=0,\quad x^6x=0.
\]

\section{The sixth-power cancellation}
\begin{theorem}\label{thm:answer}
For $S$ presented by~\eqref{eq:generators}--\eqref{eq:seven}, let
$X=\{x,y,z\}$. The sixth-power map is injective on $X$ and takes no element
of $X$ to zero. If $\omega\in\C$ is a primitive cube root of unity, then
\[
u=x+\omega y+\omega^2z\in\C X
\]
is nonzero and satisfies $u^6=0$.
\end{theorem}
\begin{proof}
The pure powers $x^6,y^6,z^6$ are distinct and nonzero by the description
of $S$ in Section~2. The coefficient of the basis vector $x$ in $u$
is $1$, so $u\ne0$.

For $v\in\C_0S$, write $[s]v$ for its coefficient at the nonzero basis
element $s\in S$. Given $a,b,c\in\C$, put $v=(ax+by+cz)^6$. Expanding
by distributivity and collecting the classes gives
\begin{align}
[x^6]v&=\sum_{j=0}^5\binom6j a^{6-j}b^j=(a+b)^6-b^6,\label{eq:cx}\\
[y^6]v&=\sum_{j=0}^5\binom6j b^{6-j}c^j=(b+c)^6-c^6,\label{eq:cy}\\
[z^6]v&=\sum_{j=0}^5\binom6j c^{6-j}a^j=(c+a)^6-a^6.\label{eq:cz}
\end{align}
These are the only possible nonzero coefficients: terms involving all
three generators are zero, and the remaining terms belong to the three
listed classes. The identities on the right are scalar binomial identities
and do not require an identity in $\C_0S$.

Take $(a,b,c)=(1,\omega,\omega^2)$. Since $1+\omega+\omega^2=0$,
\[
1+\omega=-\omega^2,\qquad \omega+\omega^2=-1,
\qquad \omega^2+1=-\omega.
\]
Using $\omega^3=1$, the three coefficients~\eqref{eq:cx}, \eqref{eq:cy}, \eqref{eq:cz} become
\[
(-\omega^2)^6-\omega^6=0,\qquad
(-1)^6-(\omega^2)^6=0,\qquad
(-\omega)^6-1=0.
\]
Thus $u^6=0$.
\end{proof}

The same proof works over any characteristic-zero field containing a
primitive cube root of unity, in particular $\mathbb Q(\omega)$.
For a coefficient-by-coefficient calculation, the binomial weights in
each list are $1,6,15,20,15,6$. The exponent of $\omega$ contributed by
$x^ay^bz^c$ is $b+2c$; its residues in each displayed list are
$0,1,2,0,1,2$ modulo three. One finds that each coefficient equals
$21(1+\omega+\omega^2)=0$.
The underlying residue-class sums are instances of the binomial identities
in Benjamin and Scott~\cite[Theorem 1.1 and p.\ 100]{BenjaminScott}.

\paragraph{A smaller example.}
Imposing, in addition to~\eqref{eq:first}--\eqref{eq:seven}, the relation $xyz=0$ gives a $49$-element semigroup with the same sixth-power
properties. Indeed, the additional zero elements are exactly the monomials
of degrees three through five divisible by $xyz$. Their number is
$1+3+6=10$; they and $0$ form an ideal, since every degree-six monomial
divisible by $xyz$ is already zero. Collapsing this ideal leaves the
three terminal elements distinct and leaves $x,y,z$ unchanged.
The preceding cancellation is therefore unchanged. Its contracted
algebra has dimension $48$ and nilpotency index seven. Denote this quotient by $\bar S$.

\section{Consequences and neighboring questions}
\begin{corollary}\label{cor:field}
Let $F$ be a field of characteristic zero, and let $T$ be either $S$ or
$\bar S$. A nonzero element of $F\{x,y,z\}\subseteq F_0T$ has sixth power
zero if and only if $F$ contains a primitive cube root of unity.
Every nonzero element in the span of any two of $x,y,z$ has nonzero sixth power.
\end{corollary}
\begin{proof}
The coefficient formulas~\eqref{eq:cx}--\eqref{eq:cz} hold over $F$ for both
presentations. If $a=0$ and all three coefficients vanish, the third formula
gives $c^6=0$ and then the second gives $b^6=0$. By cyclic symmetry,
any nonzero cancellation therefore has $abc\ne0$.
Set $r=(a+b)/b$, $s=(b+c)/c$, and $t=(c+a)/a$. Each is a sixth root of
unity different from $1$. If $F$ has no primitive cube root of unity,
the factorization
\[
q^6-1=(q-1)(q+1)(q^2+q+1)(q^2-q+1)
\]
shows that its only sixth roots of unity are $1,-1$; a root of the last
factor has its negative a primitive cube root. Thus $r=s=t=-1$.
It follows that $a=-2b$, $b=-2c$, $c=-2a$, whence $a=-8a$ and $9a=0$,
a contradiction. The converse is Theorem~\ref{thm:answer} and its quotient.
\end{proof}

In particular, the same semigroups have no such cancellation on
$\mathbb Q\{x,y,z\}$ or $\mathbb R\{x,y,z\}$, but do on
$\mathbb Q(\omega)\{x,y,z\}$. This is a property of these presentations,
not an obstruction for arbitrary semigroups over $\mathbb Q$:
Corollary~\ref{cor:rational} below gives a different example over that field.

\begin{corollary}\label{cor:multiples}
For every positive integer $n$ divisible by six and every characteristic-zero
field $F$ containing a primitive cube root of unity, there are commutative
nilpotent semigroups with zero of sizes
\[
\binom{n+2}{3}+3\quad\hbox{and}\quad \frac{3n(n-1)}2+4
\]
whose three generators have distinct nonzero $n$th powers, although a
nonzero element of their $F$-span has $n$th power zero.
Their nilpotency indices are $n+1$.
\end{corollary}
\begin{proof}
Replace six by $n$ in the three cyclic terminal lists, kill all degree-$n$
monomials involving all three generators, and kill every monomial of
degree $n+1$. As before, the three terminal classes are distinct and
all lower-degree monomials survive. The terminal coefficients of
$(ax+by+cz)^n$ are $(a+b)^n-b^n$, $(b+c)^n-c^n$, $(c+a)^n-a^n$.
At $(1,\omega,\omega^2)$ they vanish because $n$ is even and $3\mid n$.
The first count is $\sum_{d=1}^{n-1}\binom{d+2}{2}+4$.
Adding $xyz=0$ leaves $3d$ monomials in each degree $1\le d<n$, giving
the second count. Every $(n+1)$-fold product is zero and $x^n\ne0$.
\end{proof}

\begin{corollary}\label{cor:19}
The branch of Bergman's Question 19 that drops commutativity from
Question 15 also has answer no. There is even a strictly noncommutative
nilpotent example with $196$ elements at $n=6$.
\end{corollary}
\begin{proof}
The commutative example already belongs to the larger class. For the
stronger assertion let $N=\{0,p,q,r\}$ have the single nonzero product
$pq=r$. Every triple product is zero, so multiplication is associative,
and $qp=0\ne r$. The direct product $\bar S\times N$ is nilpotent and
noncommutative, has $49\cdot4=196$ elements, and contains the copy
$s\mapsto(s,0)$ of $\bar S$. Its contracted algebra contains the same
cancellation on the images of $x,y,z$.
\end{proof}

\section{Consequences of earlier identities}\label{sec:older}
Von zur Gathen and Roche~\cite[Definition 3.2, equation (3.2), p.\ 350;
Theorem 4.4(i), p.\ 357; Table 2, p.\ 359]{GathenRoche} give, for $n=35$,
the folded vector
\[
(b_0,\ldots,b_{17})=(\underbrace{1,\ldots,1}_{14},-1,-1,0,0),
\qquad \sum_{j=0}^{17}(-1)^j\binom{35}{j}b_j=0.
\]
Their folding convention for odd $n$ is $b_j=a_j-a_{n-j}$ with
$a_j\in\{0,1\}$. One unfolding takes $a_j=1$ exactly on
\[
I=\{0,1,\ldots,13,20,21\}.
\]
Thus $a_0=1$ and $a_{35}=0$: the endpoints are separated, which is
essential for distinct generator powers. The signed sum is also checked directly:
\begin{align*}
\sum_{j\in I}(-1)^j\binom{35}{j}
&=-\binom{34}{13}+\binom{35}{20}-\binom{35}{21}\\
&=\binom{34}{15}-2\binom{34}{13}=0.
\end{align*}
The first equality uses the alternating partial-sum identity; the last
uses $\binom{34}{15}/\binom{34}{13}=21\cdot20/(14\cdot15)=2$.

\begin{corollary}\label{cor:rational}
There is a $632$-element commutative nilpotent semigroup with zero
and two generators having distinct nonzero $35$th powers, while their
span in its contracted algebra over every characteristic-zero field
contains a nonzero element with $35$th power zero.
\end{corollary}
\begin{proof}
More generally, let $n\ge2$ and $I\subseteq\{0,\ldots,n\}$ satisfy
$0\in I$, $n\notin I$ and $\sum_{j\in I}(-1)^j\binom nj=0$.
Keep all positive-degree monomials in $x,y$ of degree below $n$
distinct. Kill every monomial of degree above $n$, and identify the
degree-$n$ monomials $x^{n-j}y^j$ into two distinct nonzero classes
$A,B$ according as $j\in I$ or $j\notin I$.
These relations form a semigroup congruence: every positive-degree
multiple of a terminal monomial is zero. The quotient is therefore
associative and commutative, with
\[
\sum_{d=1}^{n-1}(d+1)+3=\frac{n(n+1)}2+2
\]
elements and nilpotency index $n+1$. Its powers $x^n=A$, $y^n=B$
are distinct and nonzero. The coefficient of $A$ in $(x-y)^n$ is
the assumed signed sum; the coefficient of $B$ is zero too, since the
sum of all coefficients is $(1-1)^n=0$. The vector $x-y$ is nonzero
because $x,y$ are distinct basis elements. The displayed $I$ at $n=35$
gives the result. This is a semigroup consequence of the published
partition, rather than a claim that the 1997 paper posed or answered
Question 15.
\end{proof}

\begin{corollary}\label{cor:small19}
For every $n\ge2$, the Question-15 branch of Bergman's Question 19
has a strictly noncommutative nilpotent example of $2n+1$ elements
over every characteristic-zero field. In particular, $13$ elements
suffice at $n=6$.
\end{corollary}
\begin{proof}
The right-zero semigroup $\{e,f\}$ has multiplication $ab=b$.
Its algebra satisfies $(e-f)^2=e-f-e+f=0$, the special case of
the identity for idempotents in a common Green $\mathcal L$- or
$\mathcal R$-class in Almeida et al.~\cite[p.\ 11]{AMSV}.
For the nilpotent example, take $0$ and $(i,a)$ with $1\le i\le n$
and $a\in\{e,f\}$, setting
\[
(i,a)(j,b)=
\begin{cases}(i+j,b),&i+j\le n,\\0,&i+j>n.\end{cases}
\]
When the sum of three positive degrees is at most $n$, both bracketings
give their total degree and the rightmost label; otherwise both give zero.
This proves associativity. For $x=(1,e)$ and $y=(1,f)$, one has
$xy=(2,f)\ne(2,e)=yx$, $x^n=(n,e)\ne(n,f)=y^n$, and
$(x-y)^2=(2,e)-(2,f)-(2,e)+(2,f)=0$.
Hence $(x-y)^n=0$. All $(n+1)$-fold products vanish and $x^n\ne0$,
so the nilpotency index is $n+1$. The construction combines the cited
square-zero identity with degree truncation, as in Bergman's Example 12.
\end{proof}

\paragraph{Coverage in the posing source.}
Here $X^d$ denotes all $d$-fold products, whereas $X^{(d)}$ denotes the
set of $d$th powers. The page numbers below refer to arXiv:1206.0326v2.
The family in Corollary~\ref{cor:multiples} is not posed there; it is a
proved extension. Corollary~\ref{cor:field} is likewise a property of the
examples, not a source conjecture.

\smallskip
\noindent\begin{tabularx}{\linewidth}{@{}p{.25\linewidth}X@{}}
\toprule
Source item & Coverage \\
\midrule
Question 15, \S4, p.\ 6 & Answered: no, with $n=6$ (Theorem~\ref{thm:answer}) and over $\mathbb Q$ with $n=35$ (Corollary~\ref{cor:rational}), extending the partition of von zur Gathen and Roche (1997). \\
Question 19, \S4, p.\ 8: Question-15 branch & Answered: no (Corollaries~\ref{cor:19}, \ref{cor:small19}); the square-zero identity is reproved, with prior proof Almeida et al.\ (2007). Other branches are not covered. \\
Binomial partition and square-zero identity & Reproved; prior proofs von zur Gathen and Roche (1997), Almeida et al.\ (2007), respectively. Their semigroup consequences are proved above. \\
Question 11, \S3, p.\ 5 & Not covered as a universal question. For the displayed $X$, the product-set inequalities hold. \\
Examples 12--14, \S4, p.\ 6 & Prior constructions cited~\cite{Bergman}; Example 14 uses positive characteristic, unlike Question 15. \\
Lemma 3, \S2, p.\ 3; Corollary 10, \S3, p.\ 5 & Prior results cited~\cite{Bergman}; neither is contradicted by a span containing a vector with sixth power zero. \\
Questions 16--18, \S4, p.\ 7 & Not covered. In particular, $F X$ here fails the nonvanishing hypothesis of Question 16. \\
Conjectures 1 and 8, pp.\ 1, 5; Example 20, p.\ 8 & Not covered: the characteristic-$p$ dimension conjecture and its semigroup version are distinct, and Example 20 is nonassociative. \\
Question 23, \S6, p.\ 12 & Not covered: it asks a general annihilator-dimension inequality. \\
Corollary 4; Examples 5--6, 21--22; Lemma 9 & Not covered as separate results; these concern graded dimensions, other equality examples, or an implication between Conjectures 1 and 8. \\
\bottomrule
\end{tabularx}

\begin{corollary}\label{cor:11}
For these generating sets, the cardinalities $|X^d\setminus\{0\}|$ for $1\le d\le6$
are $(3,6,10,15,21,3)$ in $S$ and $(3,6,9,12,15,3)$ in $\bar S$.
The element $0$ belongs to $X^6$ in both examples and is excluded from these counts.
In particular both examples satisfy the inequality asked for in
Question 11 for this $X$ and $n=6$.
More generally, for the family in Corollary~\ref{cor:multiples} the counts
are $\binom{d+2}{2}$ and $3d$, respectively, for $1\le d<n$,
and both are $3$ for $d=n$, so the same compatibility holds.
\end{corollary}
\begin{proof}
For $d<6$, count all degree-$d$ monomials in $S$, or just those with at
most two generators in $\bar S$. In degree six there are exactly the
three nonzero terminal classes. Every displayed count is at least $|X|=3$.
\end{proof}

\paragraph{Verification coverage and supporting files.}
The presentations, cancellation, quotient and all six corollaries have
written proofs above. The accompanying \texttt{verify.py} exhaustively checks
both degree-six multiplication tables and the quotient map, and verifies
the cancellation by two exact calculations in $\mathbb Z[t]/(t^2+t+1)$.
The file \texttt{verify\_extensions.py} checks the product-set counts,
the four-element semigroup $N$, a finite integer-coefficient search, and
selected instances of the family in Corollary~\ref{cor:multiples}.
The family and the field criterion are general written proofs; finite
instances do not replace them.
The file \texttt{verify\_priorwork.py} checks the folded vector, its
unfolding, both signed sums, the $632$-element construction and its
$399\,424$ ordered products. It checks the truncated right-zero
constructions at $n=2,3,6$, including every associativity triple.
The general partition construction and the all-$n$ noncommutative
result have the written proofs above.

The declaration \texttt{result} in the pinned
\href{https://github.com/the-omega-institute/trureturing/blob/72243b116c47bc4fee2838e6b95fbab9802ed050/D5/S0/Certificates/Algebra/BergmanQuestion15Refutation.lean}{\texttt{BergmanQuestion15Refutation.lean}}
proves the negative answer to Question 15 with a different $59$-element
partition, displayed with its full cancellation in Appendix~\ref{app:partition}.
Writing $A=x^6$, $B=y^6$, $C=z^6$, it assigns $x^5z$ to $A$
rather than to $C$, $xy^5$ to $B$ rather than to $A$, and $yz^5$ to $C$
rather than to $B$. The simplified presentation, the $49$-element quotient
and Corollaries~\ref{cor:field}--\ref{cor:11} are covered by the written
proofs and supporting finite checks described here.
The file \texttt{verify\_original.py} checks the partition in the appendix,
all of its multiplication pairs and associativity triples, and its
cancellation by repeated multiplication and multinomial residue sums.

\paragraph{Acknowledgement.}
George M. Bergman suggested the cyclic terminal classes of $S$ as a
simplification of the construction in Appendix~\ref{app:partition}.
He also suggested the smaller example obtained by adding $xyz=0$.

\appendix
\section{A second degree-six partition}\label{app:partition}
The partition in the linked Lean statement gives a second $59$-element
semigroup, denoted $S_0$ to distinguish it from $S$.
Keep all monomials of degrees one through five distinct and identify the
six degree-six monomials in each row with the indicated element:
\begin{equation}\label{eq:original-classes}
\begin{array}{c|llllll}
A&x^6&x^3y^3&x^5y&x^5z&x^2y^4&x^4y^2\\[2pt]
B&y^6&y^3z^3&y^5z&xy^5&y^2z^4&y^4z^2\\[2pt]
C&z^6&x^3z^3&xz^5&yz^5&x^4z^2&x^2z^4.
\end{array}
\end{equation}
Send every other degree-six monomial and all monomials of degree at
least seven to zero. The rows are disjoint. Multiplying either member
of a terminal class by a positive-degree monomial gives zero, the zero
class is closed under multiplication, and every lower-degree class is a
singleton. These rules therefore define a semigroup congruence, giving
an associative commutative quotient with absorbing zero. Its cardinality
is $55+3+1=59$, its contracted complex algebra has dimension $58$, and
every product of seven elements is zero.

In $\C_0S_0$ let $w=x+\omega y+\omega^2z$, with $\omega$ a primitive
cube root of unity. Then $w\ne0$, and the generators have distinct
nonzero sixth powers $A,B,C$. A term with exponent triple $(a,b,c)$ in
$w^6$ contributes
\[
\frac{6!}{a!b!c!}\,\omega^{b+2c}x^ay^bz^c.
\]
For the rows of~\eqref{eq:original-classes}, the exponents $b+2c$ are
\[
\begin{array}{c|rrrrrr}
A&0&3&1&2&4&2\\
B&6&9&7&5&10&8\\
C&12&6&10&11&4&8.
\end{array}
\]
Each row has residues $0,0,1,2,1,2$ modulo three and weights
$1,20,6,6,15,15$. Hence each terminal coefficient is
$21(1+\omega+\omega^2)=0$, proving $w^6=0$.
This is the degree-six partition used by the declaration \texttt{result};
in its definition \texttt{positivePower}, index $5$ denotes the sixth power.

For a worked calculation, multiplication in $S_0$ gives
\[
(x^2y)(x^3)=A,\qquad (y^2z)(y^3)=B,\qquad
(xyz)(xyz)=0,\qquad Ax=0.
\]
Below degree six there are no identifications, so
\[
w^2=x^2+2\omega xy+2\omega^2xz+\omega^2y^2+2yz+\omega z^2.
\]
At degree six, $x^3y^3$ contributes $20$ to the coefficient of $A$,
while $x^2y^4$ contributes $15\omega$. In full,
\[
[A]w^6=1+20\omega^3+6\omega+6\omega^2+15\omega^4+15\omega^2
=21+21\omega+21\omega^2=0.
\]
The pure power $x^6=A$ is nonzero, but its coefficient combines with
those of mixed monomials in the same class.

\begin{thebibliography}{4}
\bibitem{Bergman}
George M. Bergman,
\emph{Thoughts on Eggert's Conjecture},
in \emph{Ring Theory and Its Applications},
Contemporary Mathematics \textbf{609} (2014), 1--17.
\href{https://doi.org/10.1090/conm/609/12101}{doi:10.1090/conm/609/12101}.
\href{https://arxiv.org/abs/1206.0326v2}{arXiv:1206.0326v2}.
\bibitem{GathenRoche}
Joachim von zur Gathen and James R.\ Roche,
\emph{Polynomials with two values},
Combinatorica \textbf{17} (1997), 345--362.
\href{https://doi.org/10.1007/BF01215917}{doi:10.1007/BF01215917}.
\bibitem{AMSV}
Jorge Almeida, Stuart Margolis, Benjamin Steinberg and Mikhail Volkov,
\emph{Representation theory of finite semigroups, semigroup radicals and
formal language theory}, 2007.
\href{https://arxiv.org/abs/math/0702400v1}{arXiv:math/0702400v1},
proof of Proposition 3.4, p.\ 11.
\bibitem{BenjaminScott}
Arthur T.\ Benjamin and Jacob N.\ Scott,
\emph{Third and fourth binomial coefficients},
The Fibonacci Quarterly \textbf{49} (2011), 99--101.
\href{https://www.fq.math.ca/Papers1/49-2/BenjaminScott.pdf}{Full text}.
\end{thebibliography}

\medskip
\noindent{\small\copyright\ 2026 Haobo Ma and Wenlin Zhang.
Licensed under \href{https://creativecommons.org/licenses/by/4.0/}{CC BY 4.0}
(cc-by-4.0).}

\end{document}
